\section{Front-Door Reliability: seguridad y disponibilidad deben cumplirse a la vez}

El servicio de identidad funciona como puerta de entrada de la organización: si deja de estar disponible, el correo, el código, el CRM, la consola de la nube y las herramientas internas pueden fallar al mismo tiempo; si concede acceso por error, un atacante puede propagarse entre aplicaciones. Por eso la reliability no puede medirse solo con el uptime, ni la security solo con los blocked attacks. El sistema necesita un SLO conjunto de availability, latency, policy correctness, state freshness y safe change.

\subsection{SLO, blast radius y safe change}

Un \term{service-level objective} (objetivo de nivel de servicio, SLO) es una meta medible que el sistema se compromete a alcanzar dentro de una ventana estadística, por ejemplo, la proporción de autenticaciones exitosas o la p99 latency. El \term{blast radius} (radio de impacto de una falla) indica el alcance de tenants, regions, applications o principals afectados por un solo error. Una plataforma de identidad multi-tenant aumenta la eficiencia al compartir recursos, pero también exige mayor isolation, canary, rollback y dependency budget.

\lecturefigure{04-front-door-reliability.png}{La identidad es una puerta de entrada con un gran radio de impacto de fallas; hay que gestionar a la vez availability, latency, correctness y change safety.}{Subtítulos locales 04:00--06:10, 14:40--16:40; redibujo conceptual.}

\begin{knowledgebox}{Lectura de la figura: las cuatro dimensiones no se sustituyen entre sí}
La availability indica si las solicitudes legítimas se completan; la latency, si la redirect, el factor y la policy chain terminan dentro de una latencia de cola aceptable; la correctness, si el allow/deny usa el principal, la policy y la session correctos; la change safety, si una regla nueva, un connector o un cambio de código pueden validarse a pequeña escala y revertirse con rapidez. Un uptime de 99.99\% no puede ocultar una autorización errónea, y una policy extremadamente estricta tampoco justifica que nadie pueda iniciar sesión.
\end{knowledgebox}

\begin{importantbox}{El error budget debe distinguir los errores de seguridad}
Un SLO común puede escribirse como
\[
\mathrm{availability}=1-\frac{N_{failed}}{N_{eligible}}.
\]
Aquí, $N_{failed}$ son las solicitudes que cumplen las condiciones estadísticas pero no tuvieron éxito, y $N_{eligible}$ son todas las eligible requests. Sin embargo, un sistema de identidad también debe rastrear por separado los false allow, los false deny, las stale entitlements y el revoke delay; estos correctness/security errors no deben quedar diluidos en un solo número global de availability.
\end{importantbox}

\subsection{Cada acceso es una context-aware policy decision}

La \term{assurance} (nivel de garantía) describe el grado de confianza del sistema en que «el sujeto actual es realmente el titular de la cuenta», y depende del authenticator, del dispositivo y del contexto. La idea común de Okta Identity Engine y de Zero Trust es que, según la aplicación de destino y el riesgo de la solicitud, la assurance requerida también cambia: una aplicación de bajo riesgo puede reutilizar la session, mientras que una operación muy sensible exige step-up authentication.

\lecturefigure{05-access-request-path.png}{La solicitud de acceso pasa por identity/context, policy, assurance y session outcome, en lugar de limitarse a verificar la contraseña.}{Okta Identity Engine policy model; NIST SP 800-207; redibujo conceptual.}

\begin{knowledgebox}{Lectura de la figura: las decisiones de policy deben poder reproducirse}
La request lleva el principal, el device/network y la target app; el context reúne profile, group, risk y session history; la policy produce allow, deny o step-up; el outcome crea la session/token y registra un audit event. Para investigar un allow anómalo, el sistema debe conservar la policy/version, los attributes clave, la risk signal, el authenticator, la decision reason y el ID de token/session, y no solo un registro de «inicio de sesión exitoso».
\end{knowledgebox}

\subsection{Resumen de la sección}

La identity reliability significa que «el sujeto correcto accede al recurso correcto con la garantía correcta, y ese acceso puede revocarse cuando cambia el riesgo». El SLO, el radio de impacto de fallas, la policy correctness y el safe rollout deben diseñarse en conjunto.
